% from ega0-22-fr.tex
\oldpage[0\textsubscript{IV-1}]{194}

  \item[(ii)] Le noyau de l'homomorphisme canonique
  $V\to V'=\mathfrak m'/\mathfrak m'^2$ est l'idéal $\mathfrak n$ de $A$
  engendré par les $\xi_i=c_{i,d_i}$ («~termes constants~» des $F_i$), et
  le conoyau de cet homomorphisme a pour base sur $K'$ les images des $T_i$~;
  en particulier, on a
  \begin{equation}
  \label{0.22.4.3.2-fr}
    \operatorname{rg}_{K'}(V')=\operatorname{rg}_K(V)+r-r',
    \tag{22.4.3.2}
  \end{equation}
  où $r'=\operatorname{rg}_K(\mathfrak n)$. Lorsque
  $\operatorname{rg}_K(V)$ est fini (autrement dit lorsque $A$ est
  artinien), pour que
  $\operatorname{rg}_{K'}(V')=\operatorname{rg}_K(V)$, il faut et il suffit
  que les $\xi_i$ soient linéairement indépendants sur $K$.
\end{enumerate}
\end{proposition}

\begin{proof}
(i) Il est clair que $A'\otimes_AK=A'/\mathfrak mA'$ est isomorphe à
$K[T_1,\ldots,T_r]/\mathfrak b'$, où $\mathfrak b'$ est l'idéal engendré
par les $r$ polynômes $T_i^{d_i}$~; c'est donc un anneau local de corps
résiduel $K$, dont l'idéal maximal est engendré par les classes des $T_i$
mod.~$\mathfrak b'$, d'où (i).

(ii) Soit $t_i$ $(1\leqslant i\leqslant r)$ la classe de $T_i$
mod.~$\mathfrak b$~; il résulte de (i) que l'idéal maximal $\mathfrak m'$ de
$A'$ est donné par
\[
  \mathfrak m'=\mathfrak mA'+\sum_{i=1}^r t_iA',
\]
d'où, compte tenu de $\mathfrak m^2=0$,
\[
  \mathfrak m'^2=\sum_{i,j}t_it_jA'+\sum_i t_i\mathfrak mA'.
\]
On en conclut que $A'/\mathfrak m'^2$ est isomorphe à l'anneau
$A[T_1,\ldots,T_r]/\mathfrak c$, où $\mathfrak c$ est l'idéal engendré par
les éléments
\[
  F_i\ (1\leqslant i\leqslant r),\qquad
  T_iT_j\ (1\leqslant i\leqslant r,\ 1\leqslant j\leqslant r),\qquad
  xT_i\ (1\leqslant i\leqslant r,\ x\in\mathfrak m).
\]
Comme $d_i\geqslant2$ et que les $T_i^2$ appartiennent à
$\mathfrak c$, les hypothèses faites entraînent que $\mathfrak c$ est aussi
engendré par les éléments
\[
  \xi_i\ (1\leqslant i\leqslant r),\qquad
  T_iT_j\ (1\leqslant i\leqslant r,\ 1\leqslant j\leqslant r),\qquad
  T_ix\ (1\leqslant i\leqslant r,\ x\in\mathfrak m).
\]
Soient
\[
  \mathfrak c_1=\sum_{i,j}T_iT_jA[T_1,\ldots,T_r],\qquad
  \mathfrak c_2=\mathfrak c_1+\sum_i\mathfrak mT_iA[T_1,\ldots,T_r],
\]
de sorte que
$\mathfrak c=\mathfrak c_2+\sum_i\xi_iA[T_1,\ldots,T_r]$. Il est clair que
$A[T_1,\ldots,T_r]/\mathfrak c_1$ est le $A$-module somme directe de $A$ et
des $At_i^{(1)}$, où $t_i^{(1)}$ est la classe de $T_i$
mod.~$\mathfrak c_1$. On en déduit que
$A[T_1,\ldots,T_r]/\mathfrak c_2$ est somme directe de $A$ et des
$Kt_i^{(2)}$, où $t_i^{(2)}$ est la classe de $T_i$
mod.~$\mathfrak c_2$. Enfin $A[T_1,\ldots,T_r]/\mathfrak c$ est somme
directe de $A/\mathfrak n$ et des $Kt_i'$, où $t_i'$ est la classe de $T_i$
mod.~$\mathfrak c$, et $\mathfrak n$ est l'idéal de $A$ engendré par les
$\xi_i$~; dans l'identification de $A'/\mathfrak m'^2$ à
$A[T_1,\ldots,T_r]/\mathfrak c$, l'image de $V$ s'identifie à
$\mathfrak m/\mathfrak n$, et ce qui précède prouve donc (ii).
\end{proof}

\begin{proposition}[22.4.4]
\label{0.22.4.4-fr}
Les hypothèses sur $A$ et les notations étant celles de (22.4.1), supposons
que $K$ soit de caractéristique $p>0$ (donc $p\cdot1\in\mathfrak m$ dans
$A$) et que les $F_i$ soient de la forme
\begin{equation}
\label{0.22.4.4.1-fr}
  F_i(T_i)=T_i^p+p\sum_{k=1}^{p-1}c_{ik}T_i^{p-k}-f_i
  \tag{22.4.4.1}
\end{equation}
avec $c_{ik}\in A$ et $f_i\in A$ pour tout $i$ et tout $k$ tel que
$1\leqslant k\leqslant p-1$~; en outre si $a_i$ est la classe de $f_i$
mod.~$\mathfrak m$, on suppose que la famille $(a_i)_{1\leqslant i\leqslant r}$
est $p$-libre sur $K^p$ («~absolument $p$-libre~»). Alors~:
\begin{enumerate}
  \item[(i)] $A'$ est un anneau local, d'idéal maximal $\mathfrak mA'$,
  dont le corps résiduel $K'$ est isomorphe à l'extension de $K$ obtenue par
  adjonction des $a_i^{1/p}$ $(1\leqslant i\leqslant r)$.

  